% Self-contained contribution. Standard AMS theorem environments and macros
% \Li, \dd, \E, \Re are supplied by the parent document.
% Bibliography keys used: herg:RadchenkoZagier, herg:DLMFhyperbolic,
% herg:DLMFbessel, herg:DLMFgamma, herg:DLMFzeta.

\section{A uniform transition for high Herglotz derivatives}
\label{herg:sec:transition}

The frozen manuscript proves a sharp asymptotic expansion for the
Herglotz derivatives when their order tends to infinity at a fixed
positive argument \cite{repoHerglotz}. Its final research paragraph explicitly leaves open
a joint limit in which the argument also varies. We resolve the
transition scale $x\asymp r$. The result has a complete expansion with
computable coefficients, an explicit finite-order remainder, and a
profile that admits arithmetic, Bessel, Gamma, and dilogarithmic
representations. The special-function identities used in the proof are
classical; the contribution here is the joint-limit theorem and its
systematic connection to this profile. No historical priority is claimed
for an isolated transformed divisor sum.

Let
\begin{equation}\label{herg:eq:F}
 F(x)=\sum_{n\geq1}\frac{\psi(nx)-\log(nx)}{n},\qquad
 D_r(x)=\frac{(-1)^{r+1}x^{r+1}}{r!}F^{(r)}(x),
 \quad x>0,\quad r\geq1.
\end{equation}
Here $F^{(r)}$ differentiates the argument $x$. Define
\begin{equation}\label{herg:eq:Bprofile}
 B(t)=\frac{1}{1-e^{-t}}-\frac1t,\qquad B(0)=\frac12,
 \qquad
 S(\lambda)=\sum_{n\geq1}\frac1{n^2}
                       B\left(\frac1{n\lambda}\right),\quad\lambda>0.
\end{equation}
The letter $B$ without an index denotes this elementary function;
$B_j$ with an integer index denotes a Bernoulli number.

\subsection{An exact probability representation and a finite remainder}

\begin{lemma}[Positive kernel and derivative bounds]
\label{herg:lem:B}
For $t>0$, $1/2<B(t)<1$ and $B'(t)>0$. For every integer $q\geq1$,
\begin{equation}\label{herg:eq:Bbound}
 \sup_{t\geq0}|B^{(q)}(t)|
 \leq \frac{2q!\zeta(q+1)}{(2\pi)^{q+1}}.
\end{equation}
\end{lemma}
\begin{proof}
The hyperbolic partial-fraction expansion
\cite[Eq.~4.36.3]{herg:DLMFhyperbolic} gives
\begin{equation}\label{herg:eq:Bpartial}
 B(t)=\frac12+2t\sum_{m\geq1}\frac1{t^2+4\pi^2m^2}
 =\frac12+\sum_{m\geq1}
       \left(\frac1{t-2\pi im}+\frac1{t+2\pi im}\right).
\end{equation}
The second sum is paired as displayed. For $q\geq1$, differentiating
gives an absolutely and uniformly convergent series on the real axis;
$|t\pm2\pi im|\geq2\pi m$ proves \eqref{herg:eq:Bbound}.
Furthermore,
\[
 B'(t)=\frac1{t^2}-\frac1{4\sinh^2(t/2)}>0,
\]
since $\sinh(t/2)>t/2$. The limits at zero and infinity are $1/2$
and $1$, respectively. The derivative bound is asserted only for
$q\geq1$; its displayed right-hand side would be divergent at $q=0$.
\end{proof}

Let $Y_r$ have the gamma density
\begin{equation}\label{herg:eq:Ydensity}
 \frac{r^{r+1}}{r!}y^r e^{-ry},\qquad y>0,
\end{equation}
and put
\begin{equation}\label{herg:eq:moments}
 \mu_m(r)=\E(Y_r-1)^m
 =\sum_{\ell=0}^m(-1)^{m-\ell}\binom m\ell
                         \frac{(r+1)_\ell}{r^\ell}.
\end{equation}
Each $\mu_m(r)$ is an exact polynomial in $r^{-1}$; it need not be
evaluated by subtracting floating-point numbers.

\begin{theorem}[Exact mixture and computable Taylor certificate]
\label{herg:thm:mixture}
For all $r\geq1$ and $\lambda>0$,
\begin{equation}\label{herg:eq:mixture}
 D_r(r\lambda)=\E S(\lambda/Y_r).
\end{equation}
Write $h_\lambda(y)=S(\lambda/y)$ for $y>0$, extended by
$h_\lambda(0)=\zeta(2)/2$. For $J\geq0$ and $q=2J+2$,
\begin{align}
 &\left|D_r(r\lambda)-
   \sum_{m=0}^{2J+1}\frac{h_\lambda^{(m)}(1)}{m!}\mu_m(r)\right|
 \label{herg:eq:certificate}\\[-2pt]
 &\hspace{18mm}\leq
 \frac{2\zeta(q+1)\zeta(q+2)}{(2\pi)^{q+1}\lambda^q}\mu_q(r).
 \notag
\end{align}
This is a finite-$r$ inequality, valid on the entire positive
$\lambda$ axis, rather than an unspecified asymptotic error.
\end{theorem}
\begin{proof}
The classical digamma integral, also used in the Herglotz theory of
Radchenko--Zagier \cite{herg:RadchenkoZagier}, is
\[
 \psi(z)-\log z=-\int_0^\infty e^{-zt}B(t)\,\dd t,
 \qquad z>0.
\]
Since $0<B<1$, the defining sum and any fixed number of its
argument derivatives converge locally uniformly. Differentiating and
substituting $u=nxt$ gives
\begin{equation}\label{herg:eq:pre-mixture}
 D_r(x)=\sum_{n\geq1}\frac1{n^2r!}
           \int_0^\infty u^r e^{-u}B\left(\frac{u}{nx}\right)\dd u.
\end{equation}
All integrands are positive and bounded after normalization by the
summable sequence $n^{-2}$. Thus every interchange is justified, and
$x=r\lambda$ proves \eqref{herg:eq:mixture}.

From \eqref{herg:eq:Bbound}, for every $q\geq1$,
\[
 \sup_{y\geq0}|h_\lambda^{(q)}(y)|
 \leq \frac{2q!\zeta(q+1)\zeta(q+2)}
                    {(2\pi)^{q+1}\lambda^q}.
\]
Taylor's formula at $y=1$, through degree $q-1$, therefore has a
global remainder bounded by the last constant times $|y-1|^q/q!$.
Taking expectations proves \eqref{herg:eq:certificate}, because $q$
is even. Formula \eqref{herg:eq:moments} follows from the elementary
gamma moments.
\end{proof}

\subsection{The complete expansion and its scope}

Define the finite polynomials $C_j(z)=\sum_m c_{j,m}z^m$ by
\begin{equation}\label{herg:eq:C}
 \sum_{j\geq0}C_j(z)w^j
 =\exp\left\{\sum_{\ell\geq1}w^\ell
          \left(\frac{z^\ell}{\ell}
                       +\frac{z^{\ell+1}}{\ell+1}\right)\right\}.
\end{equation}
Thus $\deg C_j\leq2j$. Let $\mathcal E=\lambda\,\mathrm d/\mathrm d\lambda$
and use the rising operator factorial
$\mathcal E^{\overline m}=\mathcal E(\mathcal E+1)\cdots
(\mathcal E+m-1)$, with value $1$ at $m=0$. Put
\begin{equation}\label{herg:eq:T}
 \mathcal T_j=
   \sum_m c_{j,m}(-1)^m\mathcal E^{\overline m}.
\end{equation}
This definition specifies every coefficient without an implicit
recurrence for special-function values.

\begin{theorem}[Uniform complete transition expansion]
\label{herg:thm:expansion}
For any nonnegative integers $J,u$ and any compact interval
$I\subset(0,\infty)$,
\begin{equation}\label{herg:eq:expansion}
 D_r(r\lambda)=\sum_{j=0}^J\frac{(\mathcal T_jS)(\lambda)}{r^j}
                +O_{I,J,u}(r^{-J-1}),\qquad r\longrightarrow\infty,
\end{equation}
uniformly for $\lambda\in I$, also after up to $u$ derivatives in
$\lambda$. In particular,
\begin{align}
 D_r(r\lambda)={}&S(\lambda)
  +\frac{\lambda^2S''(\lambda)}{2r}\label{herg:eq:firstcoeffs}\\
 &+\frac1{r^2}
 \left(\frac{\lambda^2S''(\lambda)}2
       +\frac{2\lambda^3S'''(\lambda)}3
       +\frac{\lambda^4S^{(4)}(\lambda)}8\right)+O_I(r^{-3}).\notag
\end{align}
\end{theorem}
\begin{proof}
The exact moment generating function is
\[
 \E e^{z(Y_r-1)}=e^{-z}(1-z/r)^{-r-1}.
\]
Its logarithm expands as the exponent in \eqref{herg:eq:C} with
$w=r^{-1}$. Therefore the coefficient of $r^{-j}$ in
$\mu_m(r)/m!$ is $[z^m]C_j(z)$. In particular
$\mu_{2J+2}(r)=O_J(r^{-J-1})$. The finite Taylor certificate proves
the complete expansion, since all central moments in it are exact
polynomials. Finally,
\[
 h_\lambda^{(m)}(1)=(-1)^m\mathcal E^{\overline m}S(\lambda),
\]
which follows by induction on $m$, proves \eqref{herg:eq:T}.

For derivative uniformity, apply the same argument to each fixed
$\partial_\lambda^v h_\lambda$. For $b\geq1$ the function
$t^bB^{(b)}(t)$ is bounded on $[0,\infty)$, by regularity at zero
and the elementary large-$t$ form $B(t)=1-t^{-1}+O(e^{-t})$.
Thus the required parameter derivatives may also be passed through
the expectation and the $n$ sum. A mixed derivative is a finite sum
of terms involving $t^bB^{(q+b)}(t)$, with $q=2J+2\geq2$.
Differentiating the paired fractions in \eqref{herg:eq:Bpartial} gives
\[
 \sup_{t\geq0}|t^b B^{(q+b)}(t)|
 \leq\frac{2(q+b)!\zeta(q+1)}{(2\pi)^{q+1}}.
\]
Indeed $t^b\leq|t\pm2\pi im|^b$. The remaining powers of
$n^{-1}$ are summable and $\lambda^{-1}$ is bounded on $I$.
This supplies the required uniform Taylor remainders after each
parameter differentiation. No uncontrolled big-$O$ term is differentiated.

For the displayed coefficients,
$C_1(z)=z+z^2/2$ and
$C_2(z)=z^2+5z^3/6+z^4/8$; substituting into
\eqref{herg:eq:T} gives \eqref{herg:eq:firstcoeffs}.
\end{proof}

\begin{corollary}[All limiting ratios and noncommuting limits]
\label{herg:cor:ratios}
The function $S$ is strictly decreasing, with
\[
 \lim_{\lambda\downarrow0}S(\lambda)=\zeta(2),\qquad
 \lim_{\lambda\to\infty}S(\lambda)=\frac{\zeta(2)}2.
\]
If $r\to\infty$ through integers and $x_r/r\to\lambda$ in the
extended interval $[0,\infty]$, then $D_r(x_r)$ tends to the
corresponding value of this continuously extended profile.
Consequently,
\[
 \lim_{x\to\infty}\lim_{r\to\infty}D_r(x)=\zeta(2),\qquad
 \lim_{r\to\infty}\lim_{x\to\infty}D_r(x)=\frac{\zeta(2)}2.
\]
The complete correction expansion is claimed on compact positive
ratio intervals; the two endpoint statements concern the leading
limit only.
\end{corollary}
\begin{proof}
Strict monotonicity follows termwise from $B'>0$. Dominated
convergence, using $1/2\leq B\leq1$, proves the limits of $S$.
Moreover $Y_r\to1$ in probability, since
$\E(Y_r-1)^2=r^{-1}+2r^{-2}$. In \eqref{herg:eq:pre-mixture}
each fixed-$n$ argument consequently tends in probability to
$1/(n\lambda)$, to infinity, or to zero, according to the ratio.
Bounded convergence in probability and the summable $n^{-2}$ bound
give the asserted three cases. At fixed $r$, letting $x\to\infty$
in \eqref{herg:eq:pre-mixture} gives $\zeta(2)/2$ directly.
\end{proof}

\section{Arithmetic continuation and zeros of the centered profile}
\label{herg:sec:profile}

Write $\sigma_{-1}(N)=\sum_{d\mid N}d^{-1}$. This section gives
exact formulas for the function that enters the joint-limit theorem,
including the complex zeros of its centered reciprocal form
$H(y)-\zeta(2)/2$. These are zeros of that centered function, not
zeros of $S$, $H$, or the original generalized Stieltjes constants.

\begin{theorem}[Arithmetic representation and meromorphic zeros]
\label{herg:thm:profilezeros}
For $\lambda>0$,
\begin{equation}\label{herg:eq:divisorprofile}
 S(\lambda)=\frac{\zeta(2)}2+
       2\lambda\sum_{N\geq1}
              \frac{\sigma_{-1}(N)}{1+4\pi^2N^2\lambda^2}.
\end{equation}
The reciprocal profile $H(y)=S(1/y)$ extends meromorphically to
the complex plane by
\begin{equation}\label{herg:eq:Hpartial}
 H(y)=\frac{\zeta(2)}2+
           2y\sum_{N\geq1}
                       \frac{\sigma_{-1}(N)}{y^2+4\pi^2N^2}.
\end{equation}
It has simple poles at $y=\pm2\pi iN$, of residue
$\sigma_{-1}(N)$. The zeros of $H(y)-\zeta(2)/2$ are precisely
the simple zero $y=0$ and the simple pairs $y=\pm i b_N$, where
there is exactly one $b_N$ in every interval
\begin{equation}\label{herg:eq:zerointervals}
 2\pi N<b_N<2\pi(N+1),\qquad N\geq1.
\end{equation}
There are no other zeros. In particular, every nonzero zero is
purely imaginary and strictly interlaces with the poles.
\end{theorem}
\begin{proof}
Insert \eqref{herg:eq:Bpartial} into \eqref{herg:eq:Bprofile}.
All terms are positive for positive arguments, so rearrangement is
permitted. Grouping the pair $(m,n)$ by $N=mn$ gives
\eqref{herg:eq:divisorprofile}. Since
$\sigma_{-1}(N)\leq1+\log N$, the series in
\eqref{herg:eq:Hpartial} converges locally uniformly away from its
displayed poles, and direct residue calculation gives their residues.

Put
\[
 T(z)=\sum_{N\geq1}\frac{\sigma_{-1}(N)}{z+4\pi^2N^2},
 \qquad H(y)-\zeta(2)/2=2yT(y^2).
\]
For nonreal $z$,
\[
 \Im T(z)=-(\Im z)\sum_{N\geq1}
             \frac{\sigma_{-1}(N)}{|z+4\pi^2N^2|^2}\ne0.
\]
For real $z>-4\pi^2$, $T(z)>0$. On each interval
$(-4\pi^2(N+1)^2,-4\pi^2N^2)$ the derivative
$T'(z)=-\sum\sigma_{-1}(j)/(z+4\pi^2j^2)^2$ is strictly
negative, while its endpoint limits are $+\infty$ and $-\infty$.
There is therefore exactly one simple zero in each interval and
none elsewhere. Taking square roots proves \eqref{herg:eq:zerointervals}.
Finally $H'(0)=2T(0)=\zeta(3)/12>0$, so the zero at the origin
is simple as well.
\end{proof}

The same representation gives positive moment and quadrature
structures. For example, for $v>0$,
\[
 \frac{S(\sqrt v)-\zeta(2)/2}{2\sqrt v}
 =\sum_{N\geq1}\frac{\sigma_{-1}(N)}{1+4\pi^2N^2v}
\]
is completely monotone in $v$: each derivative has the required
alternating sign, and the differentiated series converges locally
uniformly. This is a direct property of the displayed positive
discrete measure.

\subsection{Two complementary exact representations}

\begin{theorem}[Convergent zeta series and dual Bessel series]
\label{herg:thm:dual}
For $\lambda>1/(2\pi)$,
\begin{equation}\label{herg:eq:zetaprofile}
 S(\lambda)=\frac{\zeta(2)}2+
       \sum_{m\geq1}\frac{B_{2m}\zeta(2m+1)}{(2m)!\lambda^{2m-1}},
\end{equation}
an absolutely convergent series. For every $\lambda>0$, put
$R(\lambda)=S(\lambda)-\zeta(2)-\lambda(\log\lambda-\gamma)$.
Then
\begin{equation}\label{herg:eq:Besselprofile}
 R(\lambda)=4\lambda\sum_{N\geq1}\sigma_{-1}(N)
 \Re\left[
  \sqrt{\frac{2\pi iN}{\lambda}}\,
  K_1\left(2\sqrt{\frac{2\pi iN}{\lambda}}\right)\right],
\end{equation}
with principal square roots. The series is absolutely convergent and
locally uniformly convergent on the positive $\lambda$ axis.
\end{theorem}
\begin{proof}
The Taylor expansion of $B(t)$ at zero has radius $2\pi$.
Substitute it in \eqref{herg:eq:Bprofile}; absolute convergence
when $1/\lambda<2\pi$ gives \eqref{herg:eq:zetaprofile}.

For the second identity, Mellin inversion of $(1+u^2)^{-1}$ in
\eqref{herg:eq:divisorprofile} gives, for $1<c<2$,
\begin{equation}\label{herg:eq:Mellin}
 S(\lambda)-\frac{\zeta(2)}2
 =\frac1{2\pi i}\int_{c-i\infty}^{c+i\infty}
   \frac{\pi\zeta(s)\zeta(s+1)}{\sin(\pi s/2)}
       (2\pi)^{-s}\lambda^{1-s}\,\dd s.
\end{equation}
On this line the Dirichlet series is absolutely convergent; the
sine denominator supplies exponential vertical decay. Move the
line to $\Re s=-A$, with any fixed $A>1$. The only poles crossed
are $s=1$ and $s=0$. The negative even sine poles are cancelled
by the trivial zeros of $\zeta(s)$. Standard vertical-strip bounds
for zeta and the exponential factor justify the horizontal limits.
The residue at $1$ is $\zeta(2)/2$. At zero use
$\zeta(0)=-1/2$, $\zeta'(0)=-\frac12\log(2\pi)$, and
$\zeta(1+s)=s^{-1}+\gamma+O(s)$; the residue is
$\lambda(\log\lambda-\gamma)$.

The zeta functional equation \cite{herg:DLMFzeta}, applied twice,
transforms the remaining integral into
\begin{equation}\label{herg:eq:dualMellin}
 R(\lambda)=\frac{2\lambda}{2\pi i}\int_{(A)}
  \Gamma(w)\Gamma(w+1)\zeta(w)\zeta(w+1)
  \cos(\pi w/2)\left(\frac\lambda{2\pi}\right)^w\,\dd w.
\end{equation}
Here $\zeta(w)\zeta(w+1)=\sum\sigma_{-1}(N)N^{-w}$ is
absolutely convergent, and the product of the two gamma factors
and the cosine decays exponentially on the line. Termwise
integration is consequently justified. The standard Bessel
Mellin--Barnes formula \cite[Eq.~10.32.13]{herg:DLMFbessel} gives
\[
 \frac1{2\pi i}\int_{(A)}\Gamma(w)\Gamma(w+1)z^{-w}\,\dd w
       =2\sqrt z\,K_1(2\sqrt z),\qquad |\arg z|<\pi.
\]
Taking the average of its values at $z=2\pi iN/\lambda$ and
its conjugate produces the cosine factor and proves
\eqref{herg:eq:Besselprofile}.

For an explicit domination, set $v_N=2\sqrt{\pi N/\lambda}$.
The real-line Bessel integral implies
$|K_1(w)|\leq K_1(\Re w)\leq K_{3/2}(\Re w)$ for
$\Re w>0$. The last inequality follows from
$\cosh t\leq\cosh(3t/2)$ in that integral. Hence the absolute
value of the $N$th summand in \eqref{herg:eq:Besselprofile} is at most
\begin{equation}\label{herg:eq:sectorbound}
 2^{3/2}\pi^{3/4}\lambda^{3/4}\sigma_{-1}(N)N^{1/4}
                       e^{-v_N}(1+v_N^{-1}).
\end{equation}
This proves absolute and local uniform convergence independently
of the contour interchange.
\end{proof}

\begin{remark}[A simple numerical tail certificate]
\label{herg:rem:Besseltail}
Put $a=2\sqrt{\pi/\lambda}$ and
$C=2^{3/2}\pi^{3/4}\lambda^{3/4}$. For an integer cutoff $N\geq1$
with $a\sqrt N\geq3/2$, the absolute tail in
\eqref{herg:eq:Besselprofile} is bounded by
\begin{equation}\label{herg:eq:Besseltail}
 4C\left(1+\frac1{a\sqrt N}\right)a^{-7/2}
                  \Gamma\left(\frac72,a\sqrt N\right),
\end{equation}
where the final Gamma is the upper incomplete gamma function.
Indeed $\sigma_{-1}(n)\leq1+\log n\leq2\sqrt n$ for $n\geq1$,
and $t^{3/4}e^{-a\sqrt t}$ is decreasing for $a\sqrt t\geq3/2$.
Use the integral test and substitute $u=a\sqrt t$.
\end{remark}

\begin{corollary}[An exponentially small profile correction]
\label{herg:cor:smallratio}
As $\lambda\downarrow0$, let
\[
 A(\lambda)=2^{5/4}\pi^{3/4}\lambda^{3/4}
                       e^{-2\sqrt{\pi/\lambda}},\qquad
 \phi(\lambda)=2\sqrt{\pi/\lambda}-\frac\pi8.
\]
Then
\begin{equation}\label{herg:eq:smallratio}
 R(\lambda)=A(\lambda)
 \left[\cos\phi(\lambda)
 +\frac3{16}\sqrt{\frac\lambda{2\pi}}
       \cos\left(\phi(\lambda)+\frac\pi4\right)+O(\lambda)\right].
\end{equation}
This is additive; no division by a cosine near a zero is intended.
There are complete further powers of $\sqrt\lambda$, supplied by
the classical $K_1$ expansion. In particular $R$ takes both signs
arbitrarily close to zero.
\end{corollary}
\begin{proof}
In the $N=1$ sector, use the standard large-argument expansion
\cite[Eq.~10.40.2]{herg:DLMFbessel}
\[
 K_1(2\sqrt z)=\frac{\sqrt\pi}{2z^{1/4}}e^{-2\sqrt z}
        \left[1+\frac3{16\sqrt z}+O(z^{-1})\right]
\]
on the ray $\arg z=\pi/2$. Taking its real part gives exactly
\eqref{herg:eq:smallratio}. From \eqref{herg:eq:sectorbound},
the aggregate of the sectors $N\geq2$ is
$O(\lambda^{3/4}e^{-2\sqrt{2\pi/\lambda}})$: factor out its
$N=2$ exponential and dominate the remaining series at any fixed
positive lower value of $2\sqrt{\pi/\lambda}$. This is smaller
than every algebraic correction in the first sector. The phase
passes through arbitrarily large multiples of $\pi$ as
$\lambda\downarrow0$, proving the sign assertion.
\end{proof}

\section{Gamma and dilogarithm antiderivatives of the profile}
\label{herg:sec:primitives}

The profile admits convergent primitive identities in which the
subtractions are essential. They also provide useful test cases for
symbolic integration across the Gamma and polylogarithmic descriptions.

\begin{theorem}[Two normalized primitives]
\label{herg:thm:primitives}
For $y\geq0$ define
\begin{align}
 G(y)&=-\sum_{n\geq1}\frac1n
   \log\left[\Gamma\left(1+\frac{iy}{2\pi n}\right)
             \Gamma\left(1-\frac{iy}{2\pi n}\right)\right],
 \label{herg:eq:Gprimitive}\\
 P(y)&=\sum_{n\geq1}\left\{
     \Li_2(e^{-y/n})-\zeta(2)-\frac yn\log\frac yn
                 +\frac yn+\frac{y^2}{4n^2}\right\},
 \label{herg:eq:Pprimitive}
\end{align}
where terms at $y=0$ are taken by continuity. Each Gamma product is
strictly positive, so its logarithm is the real logarithm. Both
series converge locally uniformly, with the derivatives needed below,
and
\begin{equation}\label{herg:eq:primitivechain}
 P'(y)=G(y),\qquad G'(y)=P''(y)=H(y)-\frac{\zeta(2)}2,
 \qquad P(0)=P'(0)=G(0)=0.
\end{equation}
Equivalently,
\begin{equation}\label{herg:eq:Gammameasureprimitive}
 \int_\lambda^\infty
     \frac{S(u)-\zeta(2)/2}{u^2}\,\dd u=G(1/\lambda),
 \qquad\lambda>0.
\end{equation}
For $|y|<2\pi$ the corresponding convergent zeta-product series are
\begin{align}
 G(y)&=\sum_{m\geq1}
  \frac{B_{2m}\zeta(2m+1)}{(2m)!(2m)}y^{2m},
 \label{herg:eq:Gseries}\\
 P(y)&=\sum_{m\geq1}
  \frac{B_{2m}\zeta(2m+1)}{(2m)!(2m)(2m+1)}y^{2m+1}.
 \label{herg:eq:Pseries}
\end{align}
For example,
\[
 G(y)=\frac{\zeta(3)}{24}y^2-\frac{\zeta(5)}{2880}y^4+\cdots,
 \qquad
 P(y)=\frac{\zeta(3)}{72}y^3-\frac{\zeta(5)}{14400}y^5+\cdots.
\]
\end{theorem}
\begin{proof}
The reflection formula and the gamma recurrence
\cite[\S5.5]{herg:DLMFgamma} give, for $t\geq0$,
\[
 \Gamma\left(1+\frac{it}{2\pi}\right)
 \Gamma\left(1-\frac{it}{2\pi}\right)
     =\frac{t/2}{\sinh(t/2)}.
\]
Consequently the function
\[
 L(t)=\log\frac{\sinh(t/2)}{t/2}
      =\frac t2+\log(1-e^{-t})-\log t
\]
satisfies $L(0)=0$ and $L'(t)=B(t)-1/2$. Its primitive vanishing
at zero is
\[
 J(t)=\Li_2(e^{-t})-\zeta(2)-t\log t+t+\frac{t^2}4,
 \qquad J'(t)=L(t).
\]
Indeed $\mathrm d\Li_2(e^{-t})/\mathrm dt=\log(1-e^{-t})$.
Now $G(y)=\sum n^{-1}L(y/n)$ and $P(y)=\sum J(y/n)$.
Since $L(t)=O(t^2)$, $J(t)=O(t^3)$, and their requisite
derivatives have the corresponding bounds at zero, the series and
their stated derivatives converge uniformly on bounded positive
$y$ intervals. This proves all of \eqref{herg:eq:primitivechain}.
The substitution $y=1/u$ proves
\eqref{herg:eq:Gammameasureprimitive}. Termwise integration of the
Taylor series for $H(y)-\zeta(2)/2$ proves
\eqref{herg:eq:Gseries}--\eqref{herg:eq:Pseries}.
\end{proof}

\paragraph{Further questions arising from this theorem.}
The complete expansion above suggests a uniform matching analysis when
$\lambda=\lambda_r\downarrow0$: one must compare its derivative-dependent
remainder with the exponentially small Bessel sectors before importing
fixed-$x$ oscillatory formulas. A second problem is the asymptotic
location of the imaginary zeros $b_N$ within their certified pole
intervals; their residues are the irregular but explicit divisor weights
$\sigma_{-1}(N)$. A third is to exploit the positive Stieltjes
representation for rational upper and lower approximants to $S$, then
combine them with the exact gamma-moment remainder. Finally, logarithmic
weights in the defining Herglotz sum should yield spectral derivatives
of the divisor Dirichlet series, with zeta derivatives and Laurent
coefficients entering the transformed residue terms. These are separate
extensions; none is presumed by the compact-ratio theorem proved here.

% Suggested bibliography entries (parent should place in its bibliography):
% \bibitem{herg:RadchenkoZagier} D. Radchenko and D. Zagier,
% ``Arithmetic properties of the Herglotz function'',
% arXiv:2012.15805, https://arxiv.org/abs/2012.15805.
% \bibitem{herg:DLMFhyperbolic} NIST DLMF, Sec. 4.36,
% https://dlmf.nist.gov/4.36 (especially Eq. 4.36.3).
% \bibitem{herg:DLMFbessel} NIST DLMF, Secs. 10.32 and 10.40,
% https://dlmf.nist.gov/10.32 and https://dlmf.nist.gov/10.40.
% \bibitem{herg:DLMFgamma} NIST DLMF, Sec. 5.5,
% https://dlmf.nist.gov/5.5 (reflection and recurrence).
% \bibitem{herg:DLMFzeta} NIST DLMF, Sec. 25.4,
% https://dlmf.nist.gov/25.4 (Riemann zeta functional equation).
